% Chapter 3: The House
% Scene function: First visit inside the house. Tert catalogs the
% environment as a professional. Grief objects rendered as spatial data.
% Gordon arrives home. Turn: "I am three feet from him, in the chair
% that was hers."

\chapter{The Canvas}

The house is better than the photographs.

I come in through the back door, which is unlocked. The subject is not home. His car is not in the driveway. I checked before I crossed the street. Entering while the subject is present on a first visit is not against protocol, but it is bad practice. You want the first impression to be yours alone, without the noise of a person living in the space.

The kitchen is ordinary. Linoleum floor, which is unfortunate. Linoleum absorbs everything and returns nothing. It is the acoustic equivalent of shouting into a pillow. Whoever renovated this kitchen did not consult a haunting professional, but nobody ever does. I note the contents of the counter: a coffee maker, a single mug drying upside down on a towel, a bowl with two bananas going brown, a stack of mail that has not been opened. The refrigerator hums. It is a different hum from the office lights, but it has a similar quality of being the only sound in a room that should have more.

I walk through to the living room.

The bookshelves are everything I hoped.

Floor to ceiling, flanking the fireplace, exactly as the photographs showed. But photographs do not capture acoustics. I stop in the middle of the room and listen. Not to anything. To the room itself, to the shape of silence in it. The books absorb the top end just as I predicted. The room is quiet in a way that has texture. I rap my knuckle against the hardwood floor. Once. The sound is clean, immediate, and then it is gone. There is no ring, no sustain. I do it again near the fireplace. Same. The room eats the top and lets the bass through.

I have been in a lot of houses. Some of them were good. This one is something else.

The armchairs face slightly away from each other. I noted this in the photograph. Up close, the geometry is clearer: one has been moved about eight inches to the left, just enough to no longer face the other directly. The cushion on the unmoved chair has a deep imprint. The moved one does not. Someone sat in one chair for a long time. Then someone moved the other. Nobody has sat in the moved one since.

Above the mantel, the empty nail. The photograph on the mantel itself shows two people, but the light from the window makes it hard to see details and I do not pick it up. I am not here to learn the subject's history. I am here to learn the room.

On the bookshelf, third shelf from the top on the right side, a book lies flat on top of the other books. The same one I noticed in the survey photograph. It is a thick paperback with a cracked spine and a bookmark about two-thirds through. Everything else on these shelves has been standing upright for years. This book has been taken down recently and put back, or taken down and put back many times. I do not pick it up. It is not relevant to the work.

I walk the hallway. The grab bar is on the left wall at hip height, bolted into a stud, painted over once in the same off-white as the wall. Nobody has grabbed it recently.

At the end of the hallway, the door to the sickroom is closed.

I open it.

The hospital bed is gone. The floor has a rectangular depression where the frame sat for at least two years, based on how the hardwood has compressed. The curtains are sheer, the kind a caregiver hangs to let light in during the day without the neighbors seeing in at night. A nightstand is pushed against the wall with a glass of water on it. The water has evaporated to about a third. Nobody has touched it.

The room smells stale. Under the staleness there is something clinical I recognize: the residue of antiseptic cleaners, the kind used around medical equipment. It has soaked into the baseboards. Most people would not notice it. The subject probably does not notice it anymore. I notice it because noticing what humans have stopped noticing is the entire job.

I close the door.

The bathroom is across the hall. The raised toilet seat is still installed. The medicine cabinet is full of labeled prescription bottles, most of them expired. I read a few of the labels. They are not relevant to the work. I read them anyway.

I go back to the living room and sit down in the armchair with the imprint. I am not doing anything. I am letting the house settle around me, letting my presence become part of the room's acoustic profile so that when I begin the work, I know what the baseline sounds like with me in it.

The house creaks. Not from me. Houses creak. Old ones especially, in the afternoon, when the temperature differential between the sun-facing wall and the interior creates expansion stress in the joists. This is the sound the subject hears every day. This is the sound I need to be indistinguishable from, and then distinguishable from, in exactly the right way at exactly the right moment.

Phase 1 will be floor-level. Sub-threshold. I will use the house's own sounds and introduce variations so small that the subject's conscious mind cannot identify them but his nervous system can. He will not think something is wrong. He will feel slightly less comfortable than he did yesterday, and he will not know why.

This is the part of the work I love. Not the scaring. The calibration.

A car pulls into the driveway.

I go still. Not invisible. I am already concealed. Still in the way you go still when the subject of your attention has arrived and the observation begins.

A key in the front door. The lock turns. The door opens and closes. Footsteps on hardwood: a man who has not yet taken off his shoes because the first thing he does is put his keys on the hook and his bag on the chair by the door. I hear both of these things happen. Then the shoes come off. The footsteps become quieter and move into the kitchen. The refrigerator opens and closes. A glass of water from the tap. The footsteps come back through the hallway, past the closed sickroom door without pausing, and into the living room.

He sits down in the other armchair. The one without the imprint. The one that was moved.

He does not sit in his wife's chair.

I am three feet from him, in the chair that was hers.
